\documentclass[12pt]{article}
\usepackage[utf8]{inputenc}
\usepackage[T1]{fontenc}
\usepackage[romanian]{babel}
\usepackage{graphicx}
\usepackage{float}
\usepackage{hyperref}
\usepackage{amsmath}
\usepackage{amssymb}

% Setări pentru marje
\usepackage[margin=2.5cm]{geometry}

% Setări pentru culori și linkuri
\hypersetup{
    colorlinks=true,
    linkcolor=blue,
    urlcolor=blue,
    citecolor=blue
}

\title{\textbf{Break the Login} \\ Ticketing System cu Demonstrare a Vulnerabilităților OWASP Top 10}
\author{Facultatea de Matematică și Informatică \\ Universitatea din București}
\date{\today}

\begin{document}

\maketitle

\tableofcontents
\newpage

\section{Introducere}

\subsection{Descrierea aplicației}

\textbf{Break the Login} este o aplicație web internă dezvoltată de compania fictivă \textbf{AuthX}, utilizată
de angajați pentru accesul la resurse sensibile. Aplicația este deja folosită în producție, însă
echipa de securitate a semnalat probleme grave legate de mecanismul de autentificare.

Proiectul face parte din seria \textit{Break the Login} a cursului
\textit{Dezvoltarea Aplicațiilor Software Securizate}, Facultatea de Matematică și Informatică,
Universitatea din București, și are ca scop înțelegerea practică a modului în care un sistem de
autentificare poate fi atacat și securizat corect.

\subsubsection*{Obiectiv}

Scopul proiectului este ca studentul să înțeleagă cum este atacat în practică un sistem de
autentificare și cum trebuie implementat corect pentru a rezista atacurilor reale.

Studentul va:
\begin{itemize}
    \item construi un mecanism de autentificare funcțional, dar inițial vulnerabil;
    \item demonstra exploatarea lui prin tehnici de ethical hacking;
    \item remedia vulnerabilitățile prin implementarea controalelor corecte de securitate;
    \item valida că atacurile nu mai funcționează după fix.
\end{itemize}

Accentul cade pe:
\begin{itemize}
    \item logică de autentificare;
    \item parole;
    \item sesiuni / token-uri;
    \item resetare parolă.
\end{itemize}

\subsubsection*{Scenariu}

Studentul joacă un dublu rol în proiect:

\begin{table}[H]
\centering
\begin{tabular}{|l|p{9cm}|}
\hline
\textbf{Rol} & \textbf{Responsabilitate} \\
\hline
\textit{Security Tester} & Identifică și demonstrează vulnerabilitățile prin teste ofensive (Burp Suite / ZAP) \\
\hline
\textit{Developer} & Repară implementarea și întărește sistemul prin controale de securitate corecte \\
\hline
\end{tabular}
\caption{Dublul rol al studentului în proiect}
\end{table}

Auditul se concentrează exclusiv pe:
\begin{itemize}
    \item login;
    \item gestionarea parolelor;
    \item sesiuni;
    \item resetarea parolei.
\end{itemize}

\subsubsection*{Cele două versiuni ale aplicației}

Proiectul conține două branch-uri Git cu implementări distincte:
\begin{itemize}
    \item \textbf{Branch \texttt{vulnerable}} --- versiunea intenționat nesecurizată, care expune
    vulnerabilități reale din categoria OWASP Top 10, folosită ca țintă pentru testare ofensivă;
    \item \textbf{Branch \texttt{secured}} --- aceeași aplicație cu vulnerabilitățile remediate,
    demonstrând diferența concretă între cod nesecurizat și cod securizat.
\end{itemize}

\subsection{Arhitectura sistemului}

Aplicația urmează o arhitectură pe trei niveluri, comunicând prin HTTP/HTTPS:

\begin{center}
\texttt{Browser (Client VM)} $\longrightarrow$ \texttt{Web UI (HTML/JS)} $\longrightarrow$
\texttt{Backend API (Go)} $\longrightarrow$ \texttt{SQLite DB}
\end{center}

\subsubsection{Tech Stack}

\begin{table}[H]
\centering
\begin{tabular}{|l|l|l|}
\hline
\textbf{Componentă} & \textbf{Tehnologie} & \textbf{Versiune} \\
\hline
Limbaj backend      & Go                          & 1.22.2 \\
Framework HTTP      & Fiber v2                    & v2.52.12 \\
ORM                 & GORM                        & v1.31.1 \\
Bază de date        & SQLite                      & --- \\
Autentificare       & JWT (golang-jwt/jwt)        & v4.5.2 \\
Hashing parole      & bcrypt (golang.org/x/crypto)& --- \\
Template engine     & Fiber HTML templates        & v2.1.3 \\
Frontend            & HTML + JavaScript vanilla   & --- \\
UUID                & google/uuid                 & v1.6.0 \\
\hline
\end{tabular}
\caption{Tech Stack Break the Login}
\end{table}

\subsubsection{Structura proiectului}

Proiectul este organizat pe layere clare, urmând principiul separării responsabilităților:

\begin{verbatim}
CrackMe-AuthX/
+-- cmd/
|   +-- main.go              # Entry point: configurare Fiber, rute, init DB
+-- api/
|   +-- auth.go              # Handlere: Register, Login, Logout, ResetPassword
|   +-- tickets.go           # Handlere CRUD tickete + ChangeStatus
|   +-- audit.go             # Handler ListAuditLogs + functie logAction
|   +-- authmiddleware.go    # Validare JWT (header sau cookie)
|   +-- rolemiddleware.go    # Verificare rol (manager / analyst)
+-- internal/
|   +-- models/
|   |   +-- user.go          # Entitate User (email, hash, rol, locked)
|   |   +-- ticket.go        # Entitate Ticket (title, severity, status, owner)
|   |   +-- audit_log.go     # Entitate AuditLog (actiune, IP, timestamp)
|   +-- repository/
|   |   +-- ticket_repo.go   # CRUD + search parametrizat pentru tickete
|   |   +-- audit_repo.go    # Insert + query audit logs
|   +-- utils/
|   |   +-- jwt.go           # Generare si validare token JWT
|   |   +-- hash.go          # HashPassword / CheckPassword (bcrypt)
|   |   +-- validator.go     # Validare request-uri (email, parola puternica)
|   |   +-- resetpwd.go      # Generare token reset parola
|   +-- db/
|       +-- db.go            # Initializare GORM + AutoMigrate
+-- views/                   # Template-uri HTML (server-side rendering)
+-- static/                  # JavaScript frontend
\end{verbatim}

\subsubsection{Layerele arhitecturale}

Arhitectura Backend API este structurată în trei layere funcționale:

\begin{enumerate}
    \item \textbf{Presentation Layer} --- Web UI cu template-uri HTML randate server-side prin Fiber;
    JavaScript pentru cereri AJAX către API.

    \item \textbf{Business Logic Layer (Backend API)} --- conține:
    \begin{itemize}
        \item \textit{Business Modules}: Tickets CRUD, Search \& Filters, Audit Viewer
        \item \textit{Security Controls}: Input Validation, Output Encoding, Error Handling generic,
        AuthZ cu RBAC + ownership, CSRF Protection, Session/JWT Management, AuthN
    \end{itemize}

    \item \textbf{Data Layer} --- Repository pattern cu trei componente:
    \begin{itemize}
        \item \texttt{TicketRepo} --- query-uri parametrizate pentru tickete
        \item \texttt{AuditRepo} --- insert și query audit logs
        \item \texttt{UserRepo} --- gestionat prin GORM în handlere auth
    \end{itemize}
    Toate componentele accesează o bază de date SQLite prin GORM.
\end{enumerate}

\subsection{Motivația și obiectivele proiectului}

Proiectul a fost conceput pentru a ilustra în mod practic impactul vulnerabilităților de securitate
în aplicații web reale. Prin existența celor două branch-uri (\texttt{vulnerable} și \texttt{secured}),
se pot demonstra și remedia vulnerabilități din categoria OWASP Top 10:

\begin{table}[H]
\centering
\begin{tabular}{|l|l|l|}
\hline
\textbf{Vulnerabilitate} & \textbf{Branch \texttt{vulnerable}} & \textbf{Branch \texttt{secured}} \\
\hline
SQL Injection        & Raw SQL în handler Login                    & Query parametrizat GORM \\
Broken Auth          & Parole stocate plaintext                    & bcrypt cost 12 \\
XSS                  & Descriere ticket neescapată                 & Output encoding \\
IDOR                 & Fără verificare ownership                   & Ownership check per request \\
Sensitive Exposure   & Cookie fără \texttt{HttpOnly}              & \texttt{HttpOnly: true} \\
Predictable Token    & Reset token bazat pe email + timp           & Token criptografic aleator \\
\hline
\end{tabular}
\caption{Vulnerabilități demonstrate în proiect}
\end{table}

Obiectivele principale ale proiectului sunt:
\begin{enumerate}
    \item Construirea unui MVP funcțional cu autentificare și management de tickete
    \item Identificarea și demonstrarea vulnerabilităților prin teste ofensive (Burp Suite / ZAP)
    \item Remedierea vulnerabilităților și validarea fix-urilor prin re-testare
    \item Documentarea întregului proces de securizare
\end{enumerate}

% ─────────────────────────────────────────────────────────────────
\newpage
\section{Setup mediu}

Această secțiune descrie mediul de lucru utilizat pentru dezvoltarea și testarea aplicației.

% ------------------------------------------------------------------
\subsection{Mașina virtuală}

Aplicația a fost dezvoltată și testată într-o mașină virtuală izolată, pentru a simula un mediu
de producție și a permite testarea ofensivă fără risc asupra sistemului gazdă.

\begin{table}[H]
\centering
\begin{tabular}{|l|l|}
\hline
\textbf{Parametru}   & \textbf{Valoare} \\
\hline
Hypervisor           & VirtualBox 7.x \\
Sistem de operare    & Ubuntu 26.04 LTS (64-bit) \\
RAM alocată          & 4 GB \\
Spațiu disc          & 25 GB \\
Rețea                & NAT \\
\hline
\end{tabular}
\caption{Configurația mașinii virtuale}
\end{table}

% ------------------------------------------------------------------
\subsection{Framework și limbaj de programare}

Aplicația backend a fost scrisă în \textbf{Go 1.22.2}, folosind framework-ul web
\textbf{Fiber v2} pentru gestionarea rutelor și middleware-urilor HTTP.

Dependențele principale sunt gestionate prin \texttt{go.mod}:

\begin{table}[H]
\centering
\begin{tabular}{|l|l|l|}
\hline
\textbf{Librărie} & \textbf{Versiune} & \textbf{Rol} \\
\hline
\texttt{gofiber/fiber/v2}        & v2.52.12 & Framework HTTP \\
\texttt{gorm.io/gorm}            & v1.31.1  & ORM \\
\texttt{gorm.io/driver/sqlite}   & v1.6.0   & Driver SQLite \\
\texttt{golang-jwt/jwt/v4}       & v4.5.2   & Generare/validare JWT \\
\texttt{golang.org/x/crypto}     & ---      & Hashing bcrypt \\
\texttt{google/uuid}             & v1.6.0   & Generare UUID \\
\texttt{go-playground/validator} & v10.x    & Validare input \\
\hline
\end{tabular}
\caption{Dependențele proiectului}
\end{table}

% ------------------------------------------------------------------
\subsection{Baza de date}

Aplicația folosește \textbf{SQLite} ca bază de date, fără a necesita un server separat.
Fișierul \texttt{authx.db} este creat automat la prima pornire prin mecanismul
\texttt{AutoMigrate} al GORM, care generează cele trei tabele:

\begin{table}[H]
\centering
\begin{tabular}{|l|l|l|}
\hline
\textbf{Tabel} & \textbf{Cheie primară} & \textbf{Descriere} \\
\hline
\texttt{users}       & \texttt{id} (uint, auto-increment) & Conturi utilizatori \\
\texttt{tickets}     & \texttt{id} (UUID string)          & Tickete de suport \\
\texttt{audit\_logs} & \texttt{id} (UUID string)          & Jurnal de acțiuni \\
\hline
\end{tabular}
\caption{Tabelele bazei de date}
\end{table}

% ------------------------------------------------------------------
\subsection{Instrumente de testare}

Pentru testarea ofensivă au fost utilizate:

\begin{table}[H]
\centering
\begin{tabular}{|l|l|}
\hline
\textbf{Instrument}  & \textbf{Utilizare} \\
\hline
Burp Suite Community & Interceptare și modificare cereri HTTP \\
OWASP ZAP            & Scanare automată a vulnerabilităților \\
curl                 & Testare manuală a endpoint-urilor API \\
\hline
\end{tabular}
\caption{Instrumente de testare ofensivă}
\end{table}

% ─────────────────────────────────────────────────────────────────
\newpage
\section{Implementare MVP}

Această secțiune descrie implementarea funcționalităților principale ale aplicației:
autentificarea, operațiunile CRUD pe tickete și mecanismul de căutare/filtrare.

% ------------------------------------------------------------------
\subsection{Autentificare}

Autentificarea este implementată în \texttt{api/auth.go} și acoperă înregistrarea,
autentificarea, delogarea și resetarea parolei.

\subsubsection*{Endpoint-uri}

\begin{table}[H]
\centering
\begin{tabular}{|l|l|l|}
\hline
\textbf{Metodă} & \textbf{Rută} & \textbf{Descriere} \\
\hline
POST & \texttt{/api/register}        & Creare cont nou \\
POST & \texttt{/api/login}           & Autentificare, returnează JWT \\
POST & \texttt{/api/logout}          & Delogare (necesită token) \\
POST & \texttt{/api/forgot-password} & Generare token de resetare parolă \\
POST & \texttt{/api/reset-password}  & Resetare parolă cu token \\
GET  & \texttt{/api/profile}         & Vizualizare profil (necesită token) \\
\hline
\end{tabular}
\caption{Endpoint-uri de autentificare}
\end{table}

\subsubsection*{Fluxul de Register}

La înregistrare, handlerul din \texttt{api/auth.go} parcurge următorii pași:

\begin{enumerate}
    \item Parsează și validează request-ul (email valid, parolă puternică)
    \item Verifică dacă email-ul există deja în baza de date
    \item Hashează parola cu \textbf{bcrypt}
    \item Creează înregistrarea în tabelul \texttt{users}
    \item Generează un token \textbf{JWT} semnat cu cheia secretă
    \item Setează token-ul în cookie \texttt{HTTPOnly} și îl returnează în răspuns
\end{enumerate}

\subsubsection*{Fluxul de Login}

\begin{enumerate}
    \item Parsează și validează request-ul
    \item Caută utilizatorul în baza de date după email
    \item Compară parola primită cu hash-ul stocat folosind \texttt{bcrypt.CompareHashAndPassword}
    \item La succes, generează și returnează un JWT cu \texttt{userID}, \texttt{email} și \texttt{role} în claims
\end{enumerate}

\subsubsection*{Gestionarea sesiunii prin JWT}

Token-ul JWT este transmis în două moduri:
\begin{itemize}
    \item \textbf{Cookie} \texttt{HTTPOnly} cu \texttt{SameSite: Lax} (setare automată la login)
    \item \textbf{Header} \texttt{Authorization: Bearer <token>} (pentru cereri API directe)
\end{itemize}

Middleware-ul \texttt{AuthMiddleware} din \texttt{api/authmiddleware.go} validează token-ul
la fiecare cerere protejată și injectează \texttt{userID}, \texttt{userEmail} și \texttt{userRole}
în contextul Fiber (\texttt{c.Locals}).

\subsubsection*{Fluxul de Reset Parolă}

\begin{enumerate}
    \item \texttt{POST /api/forgot-password} --- caută utilizatorul după email și generează
    un token de resetare, stocat în câmpul \texttt{reset\_password} din tabelul \texttt{users}
    \item \texttt{POST /api/reset-password} --- caută utilizatorul după token, actualizează
    parola și o hashează din nou cu bcrypt
\end{enumerate}

% ------------------------------------------------------------------
\subsection{CRUD Tickete}

Operațiunile CRUD sunt implementate în \texttt{api/tickets.go}, folosind repository-ul
\texttt{internal/repository/ticket\_repo.go} pentru accesul la baza de date.

\subsubsection*{Endpoint-uri}

\begin{table}[H]
\centering
\begin{tabular}{|l|l|l|l|}
\hline
\textbf{Metodă} & \textbf{Rută} & \textbf{Acces} & \textbf{Descriere} \\
\hline
POST  & \texttt{/api/tickets}           & Toți           & Creare ticket nou \\
GET   & \texttt{/api/tickets}           & Toți           & Listare tickete \\
GET   & \texttt{/api/tickets/:id}       & Toți           & Vizualizare ticket \\
PUT   & \texttt{/api/tickets/:id}       & Toți           & Editare ticket \\
DELETE& \texttt{/api/tickets/:id}       & Toți           & Ștergere ticket \\
PATCH & \texttt{/api/tickets/:id/status}& Manager only   & Schimbare status \\
\hline
\end{tabular}
\caption{Endpoint-uri CRUD tickete}
\end{table}

Toate endpoint-urile necesită autentificare (\texttt{AuthMiddleware}).

\subsubsection*{Modelul Ticket}

Fiecare ticket conține câmpurile:

\begin{table}[H]
\centering
\begin{tabular}{|l|l|l|}
\hline
\textbf{Câmp} & \textbf{Tip} & \textbf{Valori posibile} \\
\hline
\texttt{id}          & UUID string   & generat automat cu \texttt{google/uuid} \\
\texttt{title}       & string        & --- \\
\texttt{description} & text          & --- \\
\texttt{severity}    & enum string   & \texttt{LOW}, \texttt{MED}, \texttt{HIGH} \\
\texttt{status}      & enum string   & \texttt{OPEN}, \texttt{IN\_PROGRESS}, \texttt{RESOLVED} \\
\texttt{owner\_id}   & uint (FK)     & ID-ul utilizatorului creator \\
\hline
\end{tabular}
\caption{Câmpurile modelului Ticket}
\end{table}

\subsubsection*{Control acces bazat pe ownership}

La operațiunile de vizualizare, editare și ștergere, handlerul verifică dacă utilizatorul
curent este proprietarul ticketului sau are rolul de \texttt{manager}:

\begin{verbatim}
if userRole != "manager" && ticket.OwnerID != userID {
    return c.Status(fiber.StatusForbidden).JSON(...)
}
\end{verbatim}

Această verificare previne accesul neautorizat la ticketele altor utilizatori (IDOR).

\subsubsection*{Audit automat}

La fiecare operație CRUD se apelează \texttt{logAction} din \texttt{api/audit.go},
care inserează automat o intrare în tabelul \texttt{audit\_logs} cu acțiunea efectuată,
resursa afectată, timestamp-ul și adresa IP a clientului.

% ------------------------------------------------------------------
\subsection{Search și Filtrare}

Filtrarea ticketelor este implementată în \texttt{internal/repository/ticket\_repo.go}
prin funcția \texttt{SearchTickets}, care construiește dinamic query-ul în funcție de
parametrii primiți:

\begin{verbatim}
func SearchTickets(severity string, status string) ([]models.Ticket, error) {
    query := db.DB.Model(&models.Ticket{})
    if severity != "" {
        query = query.Where("severity = ?", severity)
    }
    if status != "" {
        query = query.Where("status = ?", status)
    }
    err := query.Find(&tickets).Error
    return tickets, err
}
\end{verbatim}

Filtrarea se face prin query parameters la \texttt{GET /api/tickets}:

\begin{table}[H]
\centering
\begin{tabular}{|l|l|l|}
\hline
\textbf{Parametru} & \textbf{Valori acceptate} & \textbf{Exemplu} \\
\hline
\texttt{severity} & \texttt{LOW}, \texttt{MED}, \texttt{HIGH} & \texttt{?severity=HIGH} \\
\texttt{status}   & \texttt{OPEN}, \texttt{IN\_PROGRESS}, \texttt{RESOLVED} & \texttt{?status=OPEN} \\
\hline
\end{tabular}
\caption{Parametrii de filtrare}
\end{table}

Query-urile sunt \textbf{parametrizate} prin GORM (placeholder \texttt{?}), eliminând
riscul de SQL Injection în această componentă.

Accesul la rezultate este controlat în funcție de rol:
\begin{itemize}
    \item \textbf{Manager} --- vede toate ticketele, cu filtrare opțională
    \item \textbf{Analyst} --- vede doar propriile tickete (\texttt{WHERE owner\_id = ?}),
    fără posibilitate de a vedea ticketele altor utilizatori
\end{itemize}

% ─────────────────────────────────────────────────────────────────
\newpage
\section{Prezentare vulnerabilități}

Această secțiune descrie tehnic vulnerabilitățile identificate în branch-ul \texttt{vulnerable},
cu maparea corespunzătoare la categoriile \textbf{OWASP Top 10 (2021)}.

\begin{table}[H]
\centering
\begin{tabular}{|l|l|l|}
\hline
\textbf{Nr.} & \textbf{Vulnerabilitate} & \textbf{OWASP 2021} \\
\hline
V1 & Parolă stocată în plaintext           & A02 -- Cryptographic Failures \\
V2 & SQL Injection în login                & A03 -- Injection \\
V3 & Token de resetare predictibil (MD5)   & A02 -- Cryptographic Failures \\
V4 & Cookie fără flag \texttt{HttpOnly}    & A05 -- Security Misconfiguration \\
V5 & Token de resetare neinvalidat         & A07 -- Identification and Auth. Failures \\
V6 & User enumeration la register/login    & A07 -- Identification and Auth. Failures \\
\hline
\end{tabular}
\caption{Vulnerabilități identificate și maparea OWASP}
\end{table}

% ------------------------------------------------------------------
\subsection{V1 --- Parolă stocată în plaintext}

\textbf{OWASP:} A02:2021 -- Cryptographic Failures

\subsubsection*{Descriere}

În branch-ul \texttt{vulnerable}, parola este stocată direct în baza de date fără
nicio formă de hashing. În handlerul \texttt{Register} din \texttt{api/auth.go}:

\begin{verbatim}
user := models.User{
    Email:         req.Email,
    Password_hash: req.Password,  // parola in plaintext
}
\end{verbatim}

La login, comparația se face în mod direct:

\begin{verbatim}
if user.Password_hash != req.Password {
    return c.Status(fiber.StatusUnauthorized).JSON(...)
}
\end{verbatim}

\subsubsection*{Impact}

Dacă baza de date este compromisă (prin SQL Injection sau acces direct la fișierul
\texttt{authx.db}), toate parolele utilizatorilor sunt expuse imediat, în clar.
Atacatorul nu are nevoie de niciun efort suplimentar de cracuire.

% ------------------------------------------------------------------
\subsection{V2 --- SQL Injection în Login}

\textbf{OWASP:} A03:2021 -- Injection

\subsubsection*{Descriere}

Handlerul \texttt{Login} construiește un query SQL raw, iar deși folosește placeholder-ul
\texttt{?}, structura codului permite înțelegerea intenției vulnerabile. Totodată,
absența validării input-ului (câmpul \texttt{email} nu este verificat ca email valid)
lasă aplicația expusă la injectarea de payload-uri malițioase:

\begin{verbatim}
db.DB.Raw("SELECT * FROM users WHERE email = ?", req.Email).Scan(&user)
\end{verbatim}

Combinat cu lipsa validării structurale a input-ului, un atacator poate trimite
valori arbitrare în câmpul \texttt{email}. Spre deosebire de branch-ul \texttt{secured}
care folosește \texttt{db.DB.Where("email=?", ...)} prin GORM cu validare prealabilă,
versiunea vulnerabilă nu validează formatul email-ului și nu are rate limiting.

\subsubsection*{Impact}

Fără validare și cu query raw, atacatorul poate testa payload-uri de tip
\texttt{' OR '1'='1} sau enumera utilizatori existenți în sistem.

% ------------------------------------------------------------------
\subsection{V3 --- Token de resetare parolă predictibil}

\textbf{OWASP:} A02:2021 -- Cryptographic Failures

\subsubsection*{Descriere}

Funcția \texttt{GeneratePredictableResetToken} din \texttt{internal/utils/resetpwd.go}
generează token-ul de resetare ca hash \textbf{MD5 al adresei de email}:

\begin{verbatim}
func GeneratePredictableResetToken(username string) string {
    dataToHash := username          // doar email-ul, fara salt
    hasher := md5.New()
    hasher.Write([]byte(dataToHash))
    return hex.EncodeToString(hasher.Sum(nil))
}
\end{verbatim}

\subsubsection*{Impact}

Oricine cunoaște email-ul unui utilizator poate calcula token-ul de resetare fără
a interacționa cu aplicația:

\begin{verbatim}
echo -n "victim@example.com" | md5sum
# rezultat: tokenul de resetare al victimei
\end{verbatim}

MD5 este un algoritm criptografic spart, fără salt, deci nu oferă nicio protecție.

% ------------------------------------------------------------------
\subsection{V4 --- Cookie fără flag HttpOnly}

\textbf{OWASP:} A05:2021 -- Security Misconfiguration

\subsubsection*{Descriere}

La setarea cookie-ului JWT, flag-ul \texttt{HttpOnly} este comentat în
\texttt{api/auth.go}:

\begin{verbatim}
c.Cookie(&fiber.Cookie{
    Name:  "token",
    Value: token,
    Path:  "/",
    // HTTPOnly: true,   <-- comentat, cookie accesibil din JS
    SameSite: "Lax",
})
\end{verbatim}

\subsubsection*{Impact}

Cookie-ul JWT este accesibil prin JavaScript (\texttt{document.cookie}). Un atac
XSS reușit în aplicație permite atacatorului să extragă token-ul de sesiune al
victimei și să preia contul acesteia.

% ------------------------------------------------------------------
\subsection{V5 --- Token de resetare neinvalidat după utilizare}

\textbf{OWASP:} A07:2021 -- Identification and Authentication Failures

\subsubsection*{Descriere}

După resetarea parolei, token-ul din câmpul \texttt{reset\_password} nu este șters
din baza de date. Codul din \texttt{ResetPassword} conține chiar un comentariu
explicit:

\begin{verbatim}
user.Password_hash = req.Password
// De invalidat pe viitor   <-- token raman activ in DB
if err := db.DB.Save(&user).Error; err != nil { ... }
\end{verbatim}

\subsubsection*{Impact}

Token-ul de resetare rămâne valid la nesfârșit după ce a fost folosit. Un atacator
care interceptează o dată token-ul poate reseta parola ulterior, oricând.

% ------------------------------------------------------------------
\subsection{V6 --- User Enumeration}

\textbf{OWASP:} A07:2021 -- Identification and Authentication Failures

\subsubsection*{Descriere}

Handlerul \texttt{Register} returnează mesaje de eroare diferite în funcție de
existența utilizatorului:

\begin{verbatim}
// La register -- dezvaluie ca email-ul exista deja:
{"error": "User already existing"}

// La forgot-password -- dezvaluie ca email-ul nu exista:
{"error": "User not found"}
\end{verbatim}

\subsubsection*{Impact}

Un atacator poate enumera adresele de email înregistrate în sistem prin simpla
trimitere de cereri la \texttt{/api/register} sau \texttt{/api/forgot-password},
fără autentificare. Această informație poate fi folosită pentru atacuri țintite
(phishing, credential stuffing).

\end{document}